\chapter{Substituição nucleofílica}\label{ch:b1:nucleophilic-substitution}

Pegue um frasco de \cip{S}-2-bromobutano, um único \omterm{def:b1:stereochemistry-in-depth:enantiomers}{enantiômero} que gira
a luz polarizada. Aquecido com hidróxido de sódio em um \omterm{def:b1:intermolecular-forces-solvents:solvent-classes}{solvente polar}
aprótico, ele dá butan-2-ol que também é um único \omterm{def:b1:stereochemistry-in-depth:enantiomers}{enantiômero} --- mas o
outro, \cip{R}: a reação virou a molécula do avesso, como um guarda-chuva
ao vento. Deixado, em vez disso, em água morna, o mesmo brometo dá
butan-2-ol parcialmente racêmico. A mesma transformação global, um grupo
OH substituindo um átomo de Br, segue dois caminhos diferentes. Este
capítulo estabelece os dois mecanismos a partir de suas \omterm{def:b1:rate-laws:order}{leis de
velocidade}, explica sua estereoquímica e dá as regras que decidem entre
eles --- o primeiro uso completo das ferramentas dos capítulos de
cinética e de efeitos eletrônicos.

\begin{recall}
As \omterm{def:b1:rate-laws:order}{leis de velocidade}, as \omterm{def:b1:elementary-steps:elementary-step}{etapas elementares}, a \omterm{def:b1:elementary-steps:rds}{etapa determinante} e a
\omterm{def:b1:elementary-steps:qssa}{aproximação do estado estacionário} estão nos
\cref{ch:b1:rate-laws,ch:b1:elementary-steps}; as \omterm{def:b1:stereochemistry-in-depth:isomers}{configurações}, os
desenhos de Cram e a \omterm{def:b1:stereochemistry-in-depth:optical-activity}{atividade óptica}, no
\cref{ch:b1:stereochemistry-in-depth}; os \omterm{def:b1:electronic-effects:nucleophile}{nucleófilos}, os \omterm{def:b1:electronic-effects:nucleophile}{eletrófilos}, os
\omterm{def:b1:electronic-effects:nucleophile}{grupos de saída} e os \omterm{def:b1:electronic-effects:intermediates}{carbocátions}, no \cref{ch:b1:electronic-effects}; os
\omterm{def:b1:intermolecular-forces-solvents:solvent-classes}{solventes próticos} e apróticos, no
\cref{ch:b1:intermolecular-forces-solvents}.
\end{recall}

\section{A substituição e seus participantes}

\begin{definition}[Substituição nucleofílica, substrato]\label{def:b1:nucleophilic-substitution:substitution}
Uma \emph{substituição nucleofílica}\index{substituicao nucleofilica@substituição nucleofílica} é a
substituição, em um carbono saturado, de um \omterm{def:b1:electronic-effects:nucleophile}{grupo de saída} Y por um
\omterm{def:b1:electronic-effects:nucleophile}{nucleófilo} Nu: $\mathrm{Nu^-} + \mathrm{R{-}Y} \to \mathrm{R{-}Nu} +
\mathrm{Y^-}$. O composto R--Y é o \emph{substrato}\index{substrato}; o
carbono que leva Y é eletrofílico porque a ligação C--Y está polarizada
em direção a Y.
\end{definition}

Os haloalcanos são os \omterm{def:b1:nucleophilic-substitution:substitution}{substratos} típicos, com \ce{Cl-}, \ce{Br-} ou
\ce{I-} como \omterm{def:b1:electronic-effects:nucleophile}{grupos de saída}. O hidróxido dá álcoois; os alcóxidos dão
éteres; o cianeto dá nitrilas (com um carbono a mais); a amônia e as
aminas dão aminas; o iodeto se troca com os outros haletos. A água e os
álcoois, usados como solventes, também podem agir como \omterm{def:b1:electronic-effects:nucleophile}{nucleófilos}: a
substituição é então uma \omterm{def:b1:nucleophilic-substitution:sn1}{solvólise}.

\section{O mecanismo SN2}

Quando o bromometano reage com o hidróxido em água, dobrar qualquer uma
das concentrações dobra a velocidade: $v = k[\ce{CH3Br}][\ce{OH-}]$, uma
reação de ordem um em relação a cada reagente.

\begin{definition}[Mecanismo SN2, inversão de Walden]\label{def:b1:nucleophilic-substitution:sn2}
O \emph{mecanismo SN2}\index{mecanismo SN2} (\omterm{def:b1:nucleophilic-substitution:substitution}{substituição nucleofílica}
bimolecular) é uma única \omterm{def:b1:elementary-steps:elementary-step}{etapa elementar} em que o \omterm{def:b1:electronic-effects:nucleophile}{nucleófilo} ataca o
carbono pelo lado oposto ao \omterm{def:b1:electronic-effects:nucleophile}{grupo de saída} enquanto o \omterm{def:b1:electronic-effects:nucleophile}{grupo de saída}
parte. Os três outros substituintes do carbono se viram, como um
guarda-chuva que vira do avesso: é a \emph{inversão de
Walden}\index{inversao de Walden@inversão de Walden}.
\end{definition}

\begin{omfigure}
\setchemfig{atom sep=2.1em}
\schemestart
\chemfig{H@{nu}O^{-}}
\arrow{0}[,0.35]
\chemfig{@{c}C(-[:150]H_3C)(<[:210]H)(<:[:250]C_2H_5)-[@{cb}]@{br}Br}
\arrow{->}
\chemfig{[:0]HO-C(-[:30]CH_3)(<[:-30]H)(<:[:-70]C_2H_5)}
\+
\chemfig{Br^{-}}
\schemestop
\chemmove{\draw[->, thick, shorten <=2pt, shorten >=3pt](nu) .. controls +(15:5mm) and +(175:5mm) .. (c);
          \draw[->, thick, shorten <=2pt](cb) .. controls +(90:6mm) and +(90:6mm) .. (br);}

\medskip
\begin{tikzpicture}[font=\small]
  \node[draw, dashed, rounded corners, inner sep=6pt] at (0,0)
    {$\left[\vcenter{\hbox{\ce{HO}$\cdots$\chemfig{C(-[2]CH_3)(-[:-30]H)(-[:210]C_2H_5)}$\cdots$\ce{Br}}}\right]^{-}$};
  \node[below] at (0,-1.1) {estado de transição: cinco grupos em torno de C, três em um plano};
\end{tikzpicture}

{\small A reação SN2 do hidróxido com o \cip{S}-2-bromobutano. O
\omterm{def:b1:electronic-effects:nucleophile}{nucleófilo} ataca a face oposta ao bromo (\omterm{def:b1:electronic-effects:curly-arrow}{setas curvas}); no \omterm{def:b1:elementary-steps:energy-profile}{estado de
transição}, a ligação O--C se forma enquanto a ligação C--Br se rompe; o
produto é o \cip{R}-butan-2-ol: os três outros grupos passaram para o
outro lado.}
\end{omfigure}

\begin{proposition}[Lei de velocidade da SN2]\label{prop:b1:nucleophilic-substitution:sn2-law}
Para uma reação SN2, $v = k[\mathrm{RY}][\mathrm{Nu}]$.
\end{proposition}

\begin{proof}
O mecanismo é uma única \omterm{def:b1:elementary-steps:elementary-step}{etapa elementar} bimolecular; pela lei da ação
das massas para as \omterm{def:b1:elementary-steps:elementary-step}{etapas elementares} (\cref{ch:b1:elementary-steps}),
sua velocidade é proporcional ao produto das concentrações dos dois
participantes.
\end{proof}

\begin{proposition}[Estereoquímica da SN2]\label{prop:b1:nucleophilic-substitution:sn2-inversion}
Uma reação SN2 em um carbono estereogênico dá um único \omterm{def:b1:stereochemistry-in-depth:isomers}{estereoisômero},
com o arranjo invertido dos três grupos inalterados: uma reação desse
tipo é estereoespecífica.
\end{proposition}

\begin{proof}
O \omterm{def:b1:electronic-effects:nucleophile}{nucleófilo} só pode chegar ao lado livre do carbono, oposto ao \omterm{def:b1:electronic-effects:nucleophile}{grupo de
saída}; os três outros grupos passam por um plano no \omterm{def:b1:elementary-steps:energy-profile}{estado de transição}
e terminam do outro lado. A nova ligação aponta para onde a antiga não
apontava: cada molécula de \omterm{def:b1:nucleophilic-substitution:substitution}{substrato} dá o arranjo especular de seus
grupos. O descritor em geral passa de \cip{S} a \cip{R} ou vice-versa,
mas nem sempre, pois depende das prioridades do \omterm{def:b1:electronic-effects:nucleophile}{nucleófilo} e do \omterm{def:b1:electronic-effects:nucleophile}{grupo de
saída}.
\end{proof}

\begin{definition}[Reação estereoespecífica]\label{def:b1:nucleophilic-substitution:stereospecific}
Uma reação é uma \emph{reação estereoespecífica}\index{reacao estereoespecifica@reação estereoespecífica}
quando \omterm{def:b1:nucleophilic-substitution:substitution}{substratos} estereoisoméricos dão produtos estereoisomericamente
diferentes, cada \omterm{def:b1:nucleophilic-substitution:substitution}{substrato} dando o seu próprio produto.
\end{definition}

\section{O mecanismo SN1}

O 2-bromo-2-metilpropano reage com a água a uma velocidade que não
depende da concentração do \omterm{def:b1:electronic-effects:nucleophile}{nucleófilo}: $v = k[(\ce{CH3})_3\ce{CBr}]$.
Uma etapa que não envolve o \omterm{def:b1:electronic-effects:nucleophile}{nucleófilo} deve decidir a velocidade.

\begin{definition}[Mecanismo SN1, racemização, solvólise]\label{def:b1:nucleophilic-substitution:sn1}
O \emph{mecanismo SN1}\index{mecanismo SN1} (\omterm{def:b1:nucleophilic-substitution:substitution}{substituição nucleofílica}
unimolecular) tem duas etapas: a \omterm{def:b1:electronic-effects:cleavage}{heterólise} lenta da ligação C--Y, que
dá um \omterm{def:b1:electronic-effects:intermediates}{carbocátion}, e depois a adição rápida do \omterm{def:b1:electronic-effects:nucleophile}{nucleófilo} ao
\omterm{def:b1:electronic-effects:intermediates}{carbocátion}. Uma reação que transforma um único \omterm{def:b1:stereochemistry-in-depth:enantiomers}{enantiômero} em uma
mistura dos dois é uma \emph{racemização}\index{racemizacao@racemização}; uma
substituição em que o solvente é o \omterm{def:b1:electronic-effects:nucleophile}{nucleófilo} é uma
\emph{solvólise}\index{solvolise@solvólise}.
\end{definition}

\begin{omfigure}
\setchemfig{atom sep=2.1em}
\schemestart
\chemfig{C(-[2]CH_3)(-[4]H_3C)(-[6]CH_3)-[@{cb2}]@{br2}Br}
\arrow{->[lenta]}
\chemfig{C^{+}(-[2]CH_3)(-[:210]H_3C)(-[:-30]CH_3)}
\+
\chemfig{Br^{-}}
\schemestop

\smallskip
\schemestart
\chemfig{C^{+}(-[2]CH_3)(-[:210]H_3C)(-[:-30]CH_3)}
\arrow{->[rápida][\ce{H2O}]}
\chemfig{C(-[2]CH_3)(-[4]H_3C)(-[6]CH_3)-OH_2^{+}}
\arrow{->[\ce{-H+}]}
\chemfig{C(-[2]CH_3)(-[4]H_3C)(-[6]CH_3)-OH}
\schemestop
\chemmove{\draw[->, thick, shorten <=2pt](cb2) .. controls +(90:6mm) and +(90:6mm) .. (br2);}

\medskip
\begin{tikzpicture}[font=\small]
  \fill[omProp!15, draw=omProp] (0,0.2) ellipse (0.2 and 0.85);
  \fill[omProp!15, draw=omProp] (0,-0.2) ellipse (0.2 and 0.85);
  \draw[thick] (0,0) -- (-1.1,-0.2) node[left] {\ce{CH3}};
  \draw[thick] (0,0) -- (0.95,-0.45) node[right] {\ce{C2H5}};
  \draw[thick] (0,0) -- (0.4,0.5) node[above right] {H};
  \node[fill=white, inner sep=1pt] at (0,0) {\ce{C+}};
  \draw[->, thick, omDef] (-0.9,1.6) node[left] {\ce{H2O}} -- (-0.15,0.95);
  \draw[->, thick, omDef] (0.9,-1.6) node[right] {\ce{H2O}} -- (0.15,-0.95);
  \node[align=center] at (4.6,0) {ataque por qualquer face:\\ as duas configurações,\\ racêmico se nada mais\\ distinguir as faces};
\end{tikzpicture}

{\small \omterm{def:b1:nucleophilic-substitution:sn1}{Solvólise} SN1 do 2-bromo-2-metilpropano em água: ionização lenta
ao \omterm{def:b1:electronic-effects:intermediates}{carbocátion} terciário, depois captura pela água e perda de um
próton. Embaixo: um \omterm{def:b1:electronic-effects:intermediates}{carbocátion} secundário \omterm{def:b1:stereochemistry-in-depth:stereocentre}{quiral} (vindo do
2-bromobutano) é plano, e a água se adiciona a qualquer uma das faces
com igual probabilidade.}
\end{omfigure}

\begin{proposition}[Lei de velocidade da SN1]\label{prop:b1:nucleophilic-substitution:sn1-law}
Para o \omterm{def:b1:nucleophilic-substitution:sn1}{mecanismo SN1}, se o \omterm{def:b1:electronic-effects:intermediates}{carbocátion} é capturado muito mais depressa
do que volta ao \omterm{def:b1:nucleophilic-substitution:substitution}{substrato}, $v = k_1[\mathrm{RY}]$, independente do
\omterm{def:b1:electronic-effects:nucleophile}{nucleófilo}.
\end{proposition}

\begin{proof}
Etapas: $\mathrm{RY} \rightleftharpoons \mathrm{R^+} + \mathrm{Y^-}$ ($k_1$,
$k_{-1}$), $\mathrm{R^+} + \mathrm{Nu} \to \mathrm{RNu}$ ($k_2$). O
\omterm{def:b1:electronic-effects:intermediates}{carbocátion} é um \omterm{def:b1:electronic-effects:intermediates}{intermediário reativo}: pela \omterm{def:b1:elementary-steps:qssa}{aproximação do estado
estacionário} (\cref{ch:b1:elementary-steps}), $k_1[\mathrm{RY}] =
k_{-1}[\mathrm{R^+}][\mathrm{Y^-}] + k_2[\mathrm{R^+}][\mathrm{Nu}]$, logo
$v = k_2[\mathrm{R^+}][\mathrm{Nu}] = \dfrac{k_1k_2[\mathrm{RY}][\mathrm{Nu}]}
{k_{-1}[\mathrm{Y^-}] + k_2[\mathrm{Nu}]}$. Quando $k_2[\mathrm{Nu}] \gg
k_{-1}[\mathrm{Y^-}]$, isso se reduz a $k_1[\mathrm{RY}]$: a primeira
etapa é a determinante.
\end{proof}

\begin{proposition}[Estereoquímica da SN1]\label{prop:b1:nucleophilic-substitution:sn1-stereo}
Uma reação SN1 em um carbono estereogênico dá as duas \omterm{def:b1:stereochemistry-in-depth:isomers}{configurações}: a
reação não é estereoespecífica e leva à \omterm{def:b1:nucleophilic-substitution:sn1}{racemização} quando nada
distingue as duas faces do \omterm{def:b1:electronic-effects:intermediates}{carbocátion}.
\end{proposition}

\begin{proof}
O \omterm{def:b1:electronic-effects:intermediates}{carbocátion} é plano (\cref{ch:b1:electronic-effects}): suas duas faces
são imagens especulares. Se o \omterm{def:b1:electronic-effects:nucleophile}{grupo de saída} já se afastou antes da
captura, o \omterm{def:b1:electronic-effects:nucleophile}{nucleófilo} se adiciona a qualquer uma das faces com a mesma
velocidade, dando os dois \omterm{def:b1:stereochemistry-in-depth:enantiomers}{enantiômeros} em quantidades iguais. Na
prática, o íon que sai ainda protege uma face por algum tempo, e muitas
vezes se observa um ligeiro excesso de inversão.
\end{proof}

\begin{omfigure}
\begin{tikzpicture}
\begin{axis}[
  width=0.8\linewidth, height=5.0cm, xmin=0, xmax=10, ymin=-30, ymax=110,
  axis x line=bottom, axis y line=left, xlabel={coordenada de reação},
  ylabel={energia}, xtick=\empty, ytick=\empty,
  label style={font=\small}, legend style={font=\scriptsize, draw=none, at={(0.98,0.98)}, anchor=north east}]
\addplot[omDef, very thick] table[x=x, y=E] {figdata/out/bachelor-1-energy-profiles-sn2.dat};
\addplot[omProp, very thick, dashed] table[x=x, y=E] {figdata/out/bachelor-1-energy-profiles-sn1.dat};
\legend{SN2: uma etapa, SN1: duas etapas}
\node[font=\scriptsize, anchor=north] at (axis cs:5,68) {carbocátion};
\node[font=\scriptsize, anchor=north] at (axis cs:0.6,-3) {RY};
\node[font=\scriptsize, anchor=south] at (axis cs:9.6,-18) {RNu};
\end{axis}
\end{tikzpicture}

{\small \omterm{def:b1:elementary-steps:energy-profile}{Perfis de energia} esquemáticos. SN2: um único \omterm{def:b1:elementary-steps:energy-profile}{estado de
transição}, cinco grupos em torno do carbono. SN1: uma primeira barreira,
mais alta (a ionização, a \omterm{def:b1:elementary-steps:rds}{etapa determinante}), o \omterm{def:b1:electronic-effects:intermediates}{carbocátion} em um poço
raso e depois uma pequena barreira para sua captura.}
\end{omfigure}

\section{O que decide o mecanismo}

\begin{proposition}[Substrato]\label{prop:b1:nucleophilic-substitution:substrate}
Os \omterm{def:b1:nucleophilic-substitution:substitution}{substratos} metílicos e primários reagem por SN2; os terciários, por
SN1; os secundários, por qualquer um dos dois, conforme os outros
fatores. Os \omterm{def:b1:nucleophilic-substitution:substitution}{substratos} alílicos e benzílicos reagem depressa pelos dois.
\end{proposition}

\begin{proof}
A SN2 precisa de espaço atrás do carbono: cada grupo alquila no lugar de
um hidrogênio congestiona o \omterm{def:b1:elementary-steps:energy-profile}{estado de transição}, que leva cinco grupos,
e desacelera a reação; três grupos a bloqueiam. A SN1 precisa de um
\omterm{def:b1:electronic-effects:intermediates}{carbocátion} estável: a ordem de estabilidade
(\cref{prop:b1:electronic-effects:carbocation-order}) torna a ionização
rápida para os terciários, possível para os secundários e lenta demais
para os \omterm{def:b1:electronic-effects:intermediates}{carbocátions} primários e metila. Os \omterm{def:b1:electronic-effects:intermediates}{carbocátions} alílicos e
benzílicos são estabilizados por ressonância, e o mesmo sistema $\pi$
estabiliza o \omterm{def:b1:elementary-steps:energy-profile}{estado de transição} da SN2.
\end{proof}

\begin{proposition}[Grupo de saída]\label{prop:b1:nucleophilic-substitution:leaving-group}
Os dois mecanismos são tanto mais rápidos quanto melhor for o \omterm{def:b1:electronic-effects:nucleophile}{grupo de
saída}, isto é, quanto mais fraca for a base em que ele se transforma:
$\ce{I-} > \ce{Br-} > \ce{Cl-} \gg \ce{F-}$; \ce{OH-} e \ce{RO-} não
saem.
\end{proposition}

\begin{proof}
Nos dois, a ligação C--Y se rompe na \omterm{def:b1:elementary-steps:rds}{etapa determinante} ou antes dela, e
o \omterm{def:b1:elementary-steps:energy-profile}{estado de transição} leva parte da carga negativa em Y. Um grupo que
acomoda bem uma carga negativa --- a base conjugada de um \omterm{def:b1:predominant-reaction:strong-acid}{ácido forte} ---
abaixa esse \omterm{def:b1:elementary-steps:energy-profile}{estado de transição}. \ce{HI}, \ce{HBr} e \ce{HCl} são \omterm{def:b1:predominant-reaction:strong-acid}{ácidos
fortes}, \ce{HF} é um \omterm{def:b1:predominant-reaction:strong-acid}{ácido fraco} ($\mathrm{p}K_a$ 3.2), a água e os
álcoois, ácidos muito fracos: \ce{OH-} e \ce{RO-} são \omterm{def:b1:predominant-reaction:strong-acid}{bases fortes} e
maus \omterm{def:b1:electronic-effects:nucleophile}{grupos de saída}.
\end{proof}

\begin{proposition}[Solvente]\label{prop:b1:nucleophilic-substitution:solvent}
Os \omterm{def:b1:intermolecular-forces-solvents:solvent-classes}{solventes polares} próticos (água, álcoois) favorecem a SN1; os
\omterm{def:b1:intermolecular-forces-solvents:solvent-classes}{solventes polares} apróticos (propanona, dimetilsulfóxido,
dimetilformamida) favorecem a SN2 com \omterm{def:b1:electronic-effects:nucleophile}{nucleófilos} aniônicos.
\end{proposition}

\begin{proof}
A SN1 cria dois íons a partir de um \omterm{def:b1:nucleophilic-substitution:substitution}{substrato} neutro: um \omterm{def:b1:intermolecular-forces-solvents:solvent-classes}{solvente polar}
prótico estabiliza os dois, o cátion por seus \omterm{def:b1:lewis-resonance-vsepr:lewis}{pares não ligantes} e o
ânion por \omterm{def:b1:intermolecular-forces-solvents:hbond}{ligações de hidrogênio}, e abaixa a barreira de ionização. Na
SN2, a carga de um \omterm{def:b1:electronic-effects:nucleophile}{nucleófilo} aniônico se espalha no \omterm{def:b1:elementary-steps:energy-profile}{estado de
transição}; um \omterm{def:b1:intermolecular-forces-solvents:solvent-classes}{solvente prótico}, que prende fortemente o ânion pequeno
por \omterm{def:b1:intermolecular-forces-solvents:hbond}{ligações de hidrogênio}, a desacelera, enquanto um \omterm{def:b1:intermolecular-forces-solvents:solvent-classes}{solvente aprótico}
deixa o ânion livre e reativo
(\cref{ch:b1:intermolecular-forces-solvents}).
\end{proof}

\begin{method}[Escolher o mecanismo]\label{met:b1:nucleophilic-substitution:choose-mechanism}
\begin{enumerate}
  \item \omterm{def:b1:nucleophilic-substitution:substitution}{Substrato}: metila ou primário $\to$ SN2; terciário $\to$ SN1;
        secundário $\to$ continue.
  \item \omterm{def:b1:electronic-effects:nucleophile}{Nucleófilo}: forte e aniônico ($\ce{OH-}$, $\ce{RO-}$, $\ce{CN-}$,
        $\ce{I-}$) $\to$ SN2; fraco e neutro (água, álcool) $\to$ SN1.
  \item Solvente: aprótico $\to$ SN2; prótico $\to$ SN1.
  \item Preveja a estereoquímica (inversão ou \omterm{def:b1:nucleophilic-substitution:sn1}{racemização}) e verifique
        se uma \omterm{def:b1:predominant-reaction:strong-acid}{base forte} não provoca, em vez disso, uma eliminação
        (\cref{ch:b1:elimination}).
\end{enumerate}
\end{method}

\begin{omfigure}
\begin{tabular}{lll}
\hline
 & SN2 & SN1 \\
\hline
etapas & uma & duas, via um \omterm{def:b1:electronic-effects:intermediates}{carbocátion} \\
\omterm{def:b1:rate-laws:order}{lei de velocidade} & $k[\mathrm{RY}][\mathrm{Nu}]$ & $k[\mathrm{RY}]$ \\
\omterm{def:b1:nucleophilic-substitution:substitution}{substrato} & metila $>$ primário $>$ secundário & terciário $>$ secundário \\
\omterm{def:b1:electronic-effects:nucleophile}{nucleófilo} & forte, aniônico & fraco, neutro (muitas vezes o solvente) \\
solvente & polar aprótico & polar prótico \\
estereoquímica & inversão (estereoespecífica) & \omterm{def:b1:nucleophilic-substitution:sn1}{racemização} (em geral) \\
\hline
\end{tabular}

\smallskip
{\small Os dois mecanismos lado a lado.}
\end{omfigure}

\section{Exercícios}

\begin{exercise}[$\star$]\label{exo:b1:nucleophilic-substitution:1}
Escreva os produtos de: bromoetano com hidróxido de sódio; iodometano
com etóxido de sódio; 1-cloropropano com cianeto de potássio;
bromometano com amônia (primeira etapa). Identifique o \omterm{def:b1:nucleophilic-substitution:substitution}{substrato}, o
\omterm{def:b1:electronic-effects:nucleophile}{nucleófilo} e o \omterm{def:b1:electronic-effects:nucleophile}{grupo de saída}.
\end{exercise}

\begin{exercise}[$\star$]\label{exo:b1:nucleophilic-substitution:2}
Para uma reação $\mathrm{RY} + \mathrm{Nu}$ com $v = k[\mathrm{RY}][\mathrm{Nu}]$,
o que acontece com a \omterm{def:b1:rate-laws:initial-rate}{velocidade inicial} se $[\mathrm{Nu}]$ for
triplicada? E se as duas concentrações forem reduzidas à metade? Mesmas
perguntas se $v = k[\mathrm{RY}]$.
\end{exercise}

\begin{exercise}[$\star$]\label{exo:b1:nucleophilic-substitution:3}
Preveja o mecanismo (SN1 ou SN2): 1-bromobutano com iodeto de sódio em
propanona; 2-bromo-2-metilpropano em água; bromometano com hidróxido;
2-bromobutano em etanol sem base adicionada.
\end{exercise}

\begin{exercise}[$\star$]\label{exo:b1:nucleophilic-substitution:4}
Desenhe o produto da reação SN2 do cianeto com o \cip{R}-2-bromobutano
em \omterm{def:b1:stereochemistry-in-depth:representations}{representação de Cram} e dê seu descritor.
\end{exercise}

\begin{exercise}[$\star\star$]\label{exo:b1:nucleophilic-substitution:5}
Escreva o \omterm{def:b1:nucleophilic-substitution:sn1}{mecanismo SN1} completo da hidrólise do
2-bromo-2-metilpropano, com \omterm{def:b1:electronic-effects:curly-arrow}{setas curvas}, incluindo a perda do próton.
Por que a \omterm{def:b1:rate-laws:order}{lei de velocidade} é de primeira ordem, embora três espécies
participem?
\end{exercise}

\begin{exercise}[$\star\star$]\label{exo:b1:nucleophilic-substitution:6}
Explique por que o 1-bromo-2,2-dimetilpropano, um brometo primário,
reage muito lentamente por SN2 e não reage por SN1.
\end{exercise}

\begin{exercise}[$\star\star$]\label{exo:b1:nucleophilic-substitution:7}
O \cip{S}-2-iodo-octano perde sua \omterm{def:b1:stereochemistry-in-depth:optical-activity}{atividade óptica} quando mantido em
propanona com iodeto de sódio, embora nenhum composto novo se forme.
Explique e diga por que a velocidade de perda da atividade é o dobro da
velocidade de substituição.
\end{exercise}

\begin{exercise}[$\star\star$]\label{exo:b1:nucleophilic-substitution:8}
Ordene pela velocidade de SN2 com o iodeto: 2-bromo-2-metilpropano,
bromometano, 2-bromopropano, 1-bromopropano. Ordene os mesmos compostos
pela velocidade de \omterm{def:b1:nucleophilic-substitution:sn1}{solvólise} SN1.
\end{exercise}

\begin{exercise}[$\star\star$]\label{exo:b1:nucleophilic-substitution:9}
Um brometo secundário opticamente puro sofre \omterm{def:b1:nucleophilic-substitution:sn1}{solvólise}; o produto gira a
luz em \qty{12}{\percent} do valor esperado para uma inversão pura. Dê as
proporções dos dois \omterm{def:b1:stereochemistry-in-depth:enantiomers}{enantiômeros} do produto e interprete.
\end{exercise}

\begin{exercise}[$\star\star\star$]\label{exo:b1:nucleophilic-substitution:10}
Deduza a \omterm{def:b1:rate-laws:order}{lei de velocidade} geral da SN1 com a \omterm{def:b1:elementary-steps:qssa}{aproximação do estado
estacionário} e mostre que adicionar o íon do \omterm{def:b1:electronic-effects:nucleophile}{grupo de saída} \ce{Y-} no
início desacelera a reação (o \omterm{def:b1:precipitation:common-ion}{efeito do íon comum} da cinética). Em que
condição esse efeito é desprezível?
\end{exercise}

\begin{exercise}[$\star\star\star$]\label{exo:b1:nucleophilic-substitution:11}
Duas meias-vidas para a reação do bromometano com o hidróxido: com
$[\ce{OH-}]_0 = [\ce{CH3Br}]_0 = C_0$, mostre que $1/[\ce{CH3Br}] = 1/C_0 +
kt$ e dê $t_{1/2}$. Com $[\ce{OH-}]_0 = 50\,C_0$, mostre que o
decaimento é exponencial e dê $t_{1/2}$. Dados do exercício: $k =
\qty{2.0e-4}{L/(mol.s)}$, $C_0 = \qty{0.010}{mol/L}$.
\end{exercise}

\begin{exercise}[$\star\star\star$]\label{exo:b1:nucleophilic-substitution:12}
Explique por que um álcool, \ce{ROH}, não reage com o brometo de sódio,
mas reage com o ácido bromídrico concentrado dando \ce{RBr}. Escreva o
mecanismo para um álcool terciário e para um álcool primário.
\end{exercise}

\section{Problema: A solvólise do cloreto de terc-butila}

\begin{problem}[{Problema de fim de semana --- ordem e constante de
velocidade a partir de dados de condutividade, a meia-vida, o mecanismo
e sua estereoquímica, e a energia de ativação da solvólise}]
\label{pb:b1:nucleophilic-substitution:1}
O 2-cloro-2-metilpropano (cloreto de terc-butila) é dissolvido, em
pequena concentração, em uma mistura água--propanona. Sua \omterm{def:b1:nucleophilic-substitution:sn1}{solvólise}
\ce{(CH3)3CCl + H2O -> (CH3)3COH + H+ + Cl-} produz íons, e a
condutividade $\kappa$ da solução, que é proporcional à quantidade de
íons formados, é registrada. Depois de um tempo muito longo, ela chega a
$\kappa_\infty = \qty{2.000}{mS/cm}$ nas duas temperaturas. Medidas
(\unit{mS/cm}):

\begin{center}
\begin{tabular}{lcccccccc}
\hline
$t$ (min) & 0 & 5 & 10 & 15 & 20 & 30 & 45 & 60 \\
\hline
$\kappa$, \qty{25.0}{\celsius} & 0 & 0.172 & 0.329 & 0.473 & 0.605 & 0.835 & 1.110 & 1.321 \\
$\kappa$, \qty{35.0}{\celsius} & 0 & 0.507 & 0.885 & 1.168 & 1.379 & 1.654 & 1.856 & 1.940 \\
\hline
\end{tabular}
\end{center}
% data generated in figdata/bachelor-1/solvolysis.py (story data)

\medskip
\noindent\textbf{Parte I --- A ordem.}
\begin{enumerate}
  \item Por que a condutividade é proporcional à quantidade de \omterm{def:b1:nucleophilic-substitution:substitution}{substrato}
        que reagiu?
  \item Expresse $[\mathrm{RCl}]/[\mathrm{RCl}]_0$ em função de $\kappa$ e
        de $\kappa_\infty$.
  \item Escreva a lei integrada de uma reação de primeira ordem e a
        grandeza a traçar em função de $t$ para testá-la.
  \item Calcule $\ln(\kappa_\infty - \kappa)$ a \qty{25.0}{\celsius} para
        cada instante.
  \item O gráfico é uma reta? Conclua sobre a ordem.
  \item A água está em grande excesso. O que isso esconde na ordem?
\end{enumerate}

\medskip
\noindent\textbf{Parte II --- \omterm{def:b1:rate-laws:order}{Constante de velocidade} e \omterm{def:b1:rate-laws:half-life}{meia-vida}.}
\begin{enumerate}[resume]
  \item Encontre a \omterm{def:b1:rate-laws:order}{constante de velocidade} a \qty{25.0}{\celsius}, em \unit{s^{-1}}.
  \item Calcule a \omterm{def:b1:rate-laws:half-life}{meia-vida}.
  \item Em que instante \qty{90}{\percent} do \omterm{def:b1:nucleophilic-substitution:substitution}{substrato} foram consumidos?
  \item Encontre a \omterm{def:b1:rate-laws:order}{constante de velocidade} a \qty{35.0}{\celsius}.
  \item Por qual fator a velocidade aumentou com dez graus?
  \item Verifique um valor da série a \qty{35}{\celsius} com a lei.
\end{enumerate}

\medskip
\noindent\textbf{Parte III --- O mecanismo.}
\begin{enumerate}[resume]
  \item Que mecanismo o \omterm{def:b1:nucleophilic-substitution:substitution}{substrato} e o solvente sugerem?
  \item Escreva-o com \omterm{def:b1:electronic-effects:curly-arrow}{setas curvas}.
  \item Mostre que ele concorda com a ordem encontrada.
  \item Adicionar hidróxido de sódio a \qty{0.01}{mol/L} quase não altera
        a \omterm{def:b1:rate-laws:rate}{velocidade de desaparecimento} do cloreto. Por quê?
  \item O mesmo experimento com o cloreto terciário \omterm{def:b1:stereochemistry-in-depth:stereocentre}{quiral}
        3-cloro-3-metil-hexano, opticamente puro, dá um álcool de
        \omterm{def:b1:stereochemistry-in-depth:optical-activity}{atividade óptica} quase nula. Explique.
  \item Desenhe um \omterm{def:b1:elementary-steps:energy-profile}{perfil de energia} da reação e marque a \omterm{def:b1:elementary-steps:rds}{etapa
        determinante}.
  \item Por que o 1-clorobutano quase não reagiria nas mesmas condições?
\end{enumerate}

\medskip
\noindent\textbf{Parte IV --- A \omterm{def:b1:rate-laws:activation-energy}{energia de ativação}.}
\begin{enumerate}[resume]
  \item Escreva a \omterm{thm:b1:rate-laws:arrhenius}{lei de Arrhenius}.
  \item Expresse $E_a$ a partir das duas \omterm{def:b1:rate-laws:order}{constantes de velocidade}.
  \item Calcule $E_a$ ($R = \qty{8.314}{J/(K.mol)}$).
  \item O que essa energia mede, no mecanismo?
  \item Preveja a \omterm{def:b1:rate-laws:order}{constante de velocidade} a \qty{45.0}{\celsius}.
  \item Indique a \omterm{def:b1:rate-laws:activation-energy}{energia de ativação} da \omterm{def:b1:nucleophilic-substitution:sn1}{solvólise}, com dois algarismos
        significativos.
\end{enumerate}
\end{problem}
